% !TeX root = ../../Book1.tex
% 纲要标题：## 第8章　自由群与生成关系
% 纲要位置：计划纲要.md，第 333 行。

\chapter{自由群与生成关系}
\label{b1:ch:08}
\BookChapterNumber{b1:ch:08}

\begin{chapterabstract}
生成元决定元素的表达方式，关系决定哪些表达式代表同一元素。本章先以约化字构造自由群和自由积，再通过正规闭包处理群展示，并借陪集图与生成树研究自由群子群的结构。
\end{chapterabstract}

\begin{learninggoals}
\begin{itemize}
\item 能用约化字判定自由群中的等式，并从商集合的构造说明结合律。
\item 能由展示定义同态，区别“关系成立”与“展示已经证明”。
\item 理解自由积和直积的不同约束；能由陪集图及生成树构造子群自由基。
\end{itemize}
\end{learninggoals}

取两个符号 \(a,b\)，要求它们都能与各自的逆相消，却暂不规定 \(a\) 与 \(b\) 可交换。此时 \(aba^{-1}b^{-1}\) 是一个不能继续相消的字；一旦加入 \(ab=ba\) 这条关系，它就变成单位元。同一串符号在两个群中代表不同元素，差别完全来自所允许的关系。要把这种直观说法变成构造，必须先证明字的约化结果确定，再说明任意给定生成元的像何时能延拓成同态；研究群展示时，关系还会决定这个同态能否经过取商。

\ProseParagraphBreak

本章主要使用群同态、商群及生成子群。自由积与直积的比较使用\cref{b1:def:derived-series}中的交换子记号；交换化习题还用到\cref{b1:prop:derived-quotient}与\cref{b1:cor:abelianization-universal}。自由群把“先生成、后加关系”化为具体构造；\chapref{b1:ch:09}将在交换群中把同一想法写成整数关系矩阵。

\ProseParagraphBreak
\secref{b1:sec:0804}以约化字为基础，所需图和树的事实均随证明说明。

\ProseParagraphBreak
自由群与自由积的当前内容可对照 Rotman 的《An Introduction to the Theory of Groups》\cite[Chapter 11]{Rotman1995TheoryGroups}；该书第12章还讨论有限展示群字问题的不可判定性。Lyndon 与 Schupp 的《Combinatorial Group Theory》\cite[Chapter I, \S\S~1--3; Chapter II, \S\S~1--2]{LyndonSchupp2001}适合继续学习生成关系及组合群论。该书通常采用 \([h,k]=h^{-1}k^{-1}hk\)，并把群元素的作用写在对象右侧；本书的交换子与置换复合约定不同，比较共轭公式和 Schreier 代表公式时应先核对约定。

% !TeX root = ../../Book1.tex
% 纲要标题：### 8.1 自由群的构造
% 纲要位置：计划纲要.md，第 335 行。

\section{自由群的构造}
\label{b1:sec:0801}

一个生成集使每个群元素都能写成字，但不同的字可能表示同一元素。哪些等式是群公理本身强迫的，哪些来自具体群的额外限制？自由群只保留前一类等式。用 \(S\) 表示生成字母集，\(\varepsilon\) 表示空字，允许的基本变换是
\[
ss^{-1}\longleftrightarrow\varepsilon,\qquad
s^{-1}s\longleftrightarrow\varepsilon\qquad(s\in S).
\]
结合律允许重新加括号，却不改变字母的次序；我们不额外规定不同字母 \(a,b\) 满足 \(ab=ba\)，也不规定生成字母的某个正整数次幂为单位。下面将证明每个等价类有唯一约化字，因此两个约化字表示同一元素，当且仅当它们逐字相同。这就把“没有额外关系”化为可检验的等式标准，也使生成字母的像能够任意指定。

% 纲要内容（仅作写作计划，尚非教材正文）：
%
% * 约化字及其乘法
% * 自由群的泛性质
% * 基与秩
%

\subsection{约化字及其乘法}
\label{b1:subsec:0801:01}

构造尚未开始时，\(s^{-1}\) 只是另一个字母的名称。必须先规定字之间允许怎样变换，再证明所得商集合确实成为群；不能预先借用待构造群中的逆元。

\begin{definition}\label{b1:def:reduced-word}
给定集合 \(S\)，取两个不相交的符号副本 \(S\) 与 \(S^{-1}\)，把 \(s\) 与 \(s^{-1}\) 配成互逆字母，并约定 \((s^{-1})^{-1}=s\)。这些字母组成的有限序列 \(a_1\cdots a_m\) 称为字，允许长度 \(m=0\) 的空字。如果对每个 \(1\le i<m\) 都有 \(a_{i+1}\ne a_i^{-1}\)，就称这个字为\term[yuehua-zi]{约化字}[reduced word]。
\end{definition}

这里的逆号起初只是符号约定，并不预设群已经构造好。先让任意两个字用拼接相乘；拼接满足结合律，空字是单位元，但非空字在拼接运算下没有逆元。

\ProseParagraphBreak
规定插入或删除相邻互逆字母不改变字的等价类，取这些操作生成的等价关系 \(\sim\)。具体地，\(w\sim w'\) 是指存在有限链
\[
w=w_0,w_1,\ldots,w_m=w',\qquad m\ge0,
\]
其中每一步只插入或删除一对相邻互逆字母；长度为零的链允许原字保持不变。

\ProseParagraphBreak
若 \(u\sim v\)，对任意字 \(p,q\)，在上述每一步的两端加上同一前缀 \(p\) 和后缀 \(q\)，仍是允许的插入或删除。因此
\[
u\sim v\quad\Longrightarrow\quad puq\sim pvq.
\]
这正是拼接与等价关系相容的含义：改变任一因子的代表字，不会改变拼接所得的等价类。所以拼接能下降到商集合，且仍满足结合律。

\begin{lemma}\label{b1:lem:unique-reduced-word}
设 \(S\) 为集合，在字母表 \(S\sqcup S^{-1}\) 的全部有限字上，令 \(\sim\) 为由插入或删除相邻互逆字母所生成的等价关系。则每个等价类中恰有一个约化字。逐字读取原字，遇到与当前末字母互逆的字母就消去这一对，其余字母接在末尾，即可得到它。
\end{lemma}
\begin{proof}
从左到右读字，把尚未消去的字母保存在一个栈中。栈顶就是当前字的末字母：新字母若与栈顶互逆，就弹出栈顶；否则把新字母压入栈。这样保存的序列始终是约化字，并且每一步只相当于拼接一个字母后至多删除一对，所以输出与原字等价。

记向约化字读取字母 \(a\) 的状态变换为 \(T_a\)。它与 \(T_{a^{-1}}\) 互为逆操作：若原字为空，第一次接上 \(a\)，第二次删去它；若原字非空且末尾不是 \(a^{-1}\)，第一次接上 \(a\)，第二次也恰好删去它；若原末尾为 \(a^{-1}\)，第一次删去它，第二次会重新接上它，因为原约化字的倒数第二个字母不可能为 \(a\)（删除后为空时也可接上）。因此
\[
T_{a^{-1}}\circ T_a=\operatorname{id}
=T_a\circ T_{a^{-1}}.
\]

读取任何中间片段 \(aa^{-1}\) 都使保存的约化字恢复原状，后续字母的读取因而也不受影响。一次插入或删除这种片段不会改变最终输出，有限次操作也不改变。所以算法的输出只依赖等价类，不依赖所选的插入、删除链。约化字本身是归约过程的不动输出，任何约化代表都必须等于这个唯一输出；无论先消去哪一对，只要最后得到约化字，结果就相同。
\end{proof}

对任意集合 \(S\)，上述过程都给出有限字的归约规则。若字母可以有效编码，且能判定两个字母是否互逆，\cref{b1:alg:free-word-reduction}就把它写成可执行的算法；固定有限字母表满足这些条件。算法读完输入字后必定停机，其正确性正由上面的引理保证。

\begin{algorithm}[htbp]
\caption{自由字约化}
\label{b1:alg:free-word-reduction}
\begin{algorithmsteps}{字母表 \(S\sqcup S^{-1}\) 上的有限字 \(w=a_1\cdots a_m\)，其中字母可有效编码并可判定两个字母是否互逆。}{与 \(w\) 等价的唯一约化字。}
\State 置空字 \(L\)。
\For{\(j=1,\ldots,m\)}
\If{\(L\) 非空，且 \(a_j\) 与 \(L\) 的末字母互逆}
\State 删除 \(L\) 的末字母。
\Else
\State 将 \(a_j\) 接在 \(L\) 的末尾。
\EndIf
\EndFor
\State \Return \(L\)。
\end{algorithmsteps}
\end{algorithm}

\begin{definition}[自由群的约化字构造]\label{b1:def:free-group-construction}
给定集合 \(S\)，在字母表 \(S\sqcup S^{-1}\) 的全部有限字上，令 \(\sim\) 为由插入或删除相邻互逆字母所生成的等价关系。令 \(F(S)\) 为这些字关于 \(\sim\) 的等价类集合，以 \([w]\) 表示字 \(w\) 的等价类，并配备由拼接诱导的乘法
\[
[u]\cdot[v]=[uv].
\]
下段验证这一运算使 \(F(S)\) 成为群。经此构造得到的 \(F(S)\) 称为以 \(S\) 为自由生成字母集的\term[ziyou-qun]{自由群}[free group]。用每个类的唯一约化字表示其元素时，乘法就是“先拼接，再约化”。
\end{definition}

这一乘法的良定义性与结合律已由拼接同等价关系的相容性得到。空字所在的类为单位；字 \(a_1\cdots a_m\) 的逆由 \(a_m^{-1}\cdots a_1^{-1}\) 表示，两个次序的拼接均可从中间向外成对消去为 \(\varepsilon\)。因此该构造确实给出一个群。
\symidx{\(F(S)\)：以 \(S\) 为基的自由群}

\ProseParagraphBreak
取两个不同字母 \(a,b\)。在 \(F\bigl(\{a,b\}\bigr)\) 中有
\[
ab^{-1}bba^{-1}\rightsquigarrow aba^{-1},\qquad
[ab^{-1}bba^{-1}]=[aba^{-1}].
\]
左右形式字的长度分别为五和三，作为字母序列并不相同，却在 \(F\bigl(\{a,b\}\bigr)\) 中代表同一元素。另一方面，\(aba^{-1}b^{-1}\) 已约化且非空，所以它所在的类不同于空字的类，代表的元素不是单位。这给出一条直接可判定的等式标准：先各自约化，再逐字比较。

\subsection{自由群的泛性质}
\label{b1:subsec:0801:02}

自由群中的字母只受互逆相消的规定约束。因此，任给目标群 \(G\)，每个 \(s\in S\) 都可以彼此独立地送到 \(G\) 的任意元素，逆字母的像再由这些取值确定。无需要求不同字母的像交换，也无需规定它们满足某个有限阶关系；下述泛性质中的“任意集合映射”正是这个意思。

\begin{theorem}\label{b1:thm:free-group-universal}
设 \(S\) 为集合，\(F(S)\) 为\cref{b1:def:free-group-construction}构造的自由群，\(\iota:S\to F(S)\) 将字母映为长度为一的字。对任意群 \(G\) 及集合映射 \(f:S\to G\)，存在唯一群同态 \(\widetilde f:F(S)\to G\) 满足 \(\widetilde f\circ\iota=f\)。
\end{theorem}
\begin{proof}
把 \(s^{-1}\) 映为 \(f(s)^{-1}\)，再把每个字映为其字母像在 \(G\) 中的乘积。插入或删除相邻互逆字母不改变乘积，所以映射在等价类上良定义。拼接对应乘法，因而所得 \(\widetilde f\) 是同态。任意满足条件的同态在生成字母及其逆上的取值均已确定，所以在每个字上的取值也确定，给出唯一性。
\end{proof}

唯一延拓的关系如下；图按底集合理解，\(\iota,f\) 是集合映射，虚线表示唯一的群同态：
\[
\begin{tikzcd}[column sep=3em,row sep=2.5em]
S\arrow[r,"\iota"]\arrow[dr,"f"']&F(S)\arrow[d,dashed,"\exists!\,\widetilde f"]\\
&G
\end{tikzcd}
\]
这称为自由群的泛性质。这里任意的是集合映射 \(f\)，它不需要先满足交换性、有限阶或其他关系。例如指定 \(a\mapsto(12)\)、\(b\mapsto(123)\) 就得到 \(F\bigl(\{a,b\}\bigr)\to S_3\) 的满同态。该同态的核不平凡，因为 \(a^2\) 在自由群中非单位，却映为单位。

\begin{proposition}\label{b1:prop:free-universal-unique}
设 \(S\) 为集合，群 \(F\) 配有映射 \(j:S\to F\)。假设对任意群 \(G\) 和集合映射 \(f:S\to G\)，都存在唯一群同态 \(v:F\to G\) 满足 \(v\circ j=f\)。则有唯一同构 \(F(S)\to F\) 把每个 \(\iota(s)\) 映为 \(j(s)\)，其中 \(F(S)\) 与 \(\iota:S\to F(S)\) 如\cref{b1:thm:free-group-universal}。
\end{proposition}
\begin{proof}
沿用\cref{b1:thm:universal-uniqueness}中比较两个自由幺半群的方法：先互相取作目标，构造两个方向的同态，再用延拓唯一性判定复合。将\cref{b1:thm:free-group-universal}用于映射 \(j:S\to F\)，得到唯一同态 \(u:F(S)\to F\)，满足 \(u\iota=j\)。再把假设中 \(F\) 的泛性质用于 \(\iota:S\to F(S)\)，得到唯一同态 \(v:F\to F(S)\)，满足 \(vj=\iota\)。

对于第一个复合，有
\[
(vu)\circ\iota=\iota=\operatorname{id}_{F(S)}\circ\iota.
\]
因此 \(vu\) 与恒等同态都是同一集合映射 \(\iota:S\to F(S)\) 的同态延拓。这里使用 \(F(S)\) 的泛性质中的唯一性，得到 \(vu=\operatorname{id}_{F(S)}\)。对于反向复合，同样有
\[
(uv)\circ j=j=\operatorname{id}_F\circ j.
\]
这一次，\(uv\) 与恒等同态都延拓 \(j:S\to F\)，使用的是假设中 \(F\) 的泛性质的唯一性，故 \(uv=\operatorname{id}_F\)。因此 \(u\) 为同构。任意满足所述条件的同构也延拓 \(j\)，仍由第一个泛性质的唯一性等于 \(u\)。
\end{proof}

“唯一”必须连同生成映射一起理解；若忘掉生成映射，自由群通常有很多自同构。泛性质确定的是带有指定生成字母的对象。

\subsection{基与秩}
\label{b1:subsec:0801:03}

\begin{definition}\label{b1:def:free-basis-rank}
设 \(F\) 为群，\(B\subseteq F\)。若每个 \(g\in F\) 都能唯一写成由 \(B\) 的元素及其逆组成的有限约化字，则称 \(B\) 为 \(F\) 的\term[ziyou-ji-basis]{自由基}[free basis]。这里“约化”指相邻位置不是同一个基符号的一对互逆字母，空字表示单位元；唯一性比较的是带正、逆标记的字母序列。具有这样的基的群称为自由群，基的基数称为\term[ziyou-qun-de-zhi]{自由群的秩}[rank of a free group]。\cref{b1:prop:free-rank-invariant}将证明此基数不依赖自由基的选择，因此秩是群本身的不变量。
\end{definition}

等价地，从符号集合 \(B\) 构造的自由群 \(F(B)\) 到 \(F\) 的同态（把每个符号送到 \(F\) 中的同名元素）是同构：存在字表示给出满射性，唯一约化字表示给出单射性。因此，前面构造的 \(F(S)\) 确实以单字母为自由基，而任意具有自由基的群也同构于这一构造。

\ProseParagraphBreak
一般生成集未必是自由基。\(F\bigl(\{a,b\}\bigr)\) 中的 \(\{a,b,ab\}\) 是生成集，却满足三个指定元素之间的非平凡关系；它不能作为三个自由字母。相反，\(\{a,ab\}\) 是自由基。由\cref{b1:thm:free-group-universal}，以下指定分别延拓为 \(F\bigl(\{a,b\}\bigr)\) 到自身的同态 \(\phi,\psi\)：
\[
\begin{aligned}
\phi(a)&=a,&\phi(b)&=ab,\\
\psi(a)&=a,&\psi(b)&=a^{-1}b.
\end{aligned}
\]
两个方向的复合可逐项计算：
\[
\begin{aligned}
(\psi\phi)(a)&=a,&
(\psi\phi)(b)&=\psi(a)\psi(b)=a(a^{-1}b)=b,\\
(\phi\psi)(a)&=a,&
(\phi\psi)(b)&=\phi(a)^{-1}\phi(b)=a^{-1}(ab)=b.
\end{aligned}
\]
由延拓唯一性，两个复合均为恒等同态，所以 \(\phi\) 是自同构。它把原自由基 \(\{a,b\}\) 送到 \(\{a,ab\}\)；对后者的字表达式施以 \(\phi^{-1}\)，就回到原基的唯一约化字表示，故后者也是自由基。

\begin{proposition}\label{b1:prop:free-rank-invariant}
设 \(F\) 为自由群。则 \(F\) 的任意两组自由基具有相同基数。
\end{proposition}
\begin{proof}
设 \(S,T\) 为 \(F\) 的两组自由基。由前面的同构 \(F\cong F(S)\) 及\cref{b1:thm:free-group-universal}，\(F\) 到 \(C_2\) 的同态与集合映射 \(S\to C_2\) 一一对应。因此若 \(|S|=r<\infty\)，同态数恰为 \(2^r\)，因为每个生成字母的像有两个独立选择。若两基都有限，记 \(|T|=s\)，则同一个群的同态数给出 \(2^r=2^s\)，有限整数的指数计数遂给出 \(r=s\)。

还需排除有限与无限的混合情况：若 \(F(S)\) 可由有限多个字生成，这些字总共只涉及 \(S\) 的有限子集 \(S_0\)。它们的乘积及约化仍只涉及 \(S_0\)，由约化字的唯一性，不可能产生 \(S\setminus S_0\) 中的单字母。因此必须有 \(S=S_0\)。这一论证适用于任意自由基，所以有限生成自由群的每个自由基都有限。

对于无限基数，即使仍把同态数写成 \(2^{|S|}\) 与 \(2^{|T|}\)，上面的有限整数计数也没有给出从指数值相等反推基数相等的论证。因此改用每个基元素只涉及有限个字母这一事实，直接比较两组基。

若 \(S,T\) 都无限，把每个 \(t\in T\) 写成有限个 \(S\) 字母的字，记所涉及的基字母组成的有限子集为 \(S_t\)。因为 \(T\) 生成全群，有 \(\bigcup_{t\in T}S_t=S\)；否则这些字的乘积、取逆及约化仍无法产生缺失的单字母。此处采用选择公理，并使用附录中的\cref{b1:prop:finite-support-cardinality}：由无限集合 \(T\) 标号的有限集之并，基数至多为 \(|T|\)，且 \(|T\times\mathbb N|=|T|\)。因而
\[
|S|\le |T|\cdot\aleph_0=|T|.
\]
交换 \(S,T\) 得到 \(|T|\le|S|\)，再由\cref{b1:prop:cantor-bernstein}中“互有单射则等势”的结论得到 \(|S|=|T|\)。
\end{proof}

记秩为 \(r\) 的自由群为 \(F_r\)。\(F_0\) 为单位群，\(F_1\cong\mathbb Z\)；当 \(r\ge2\) 时，两个不同基元素的交换子是非空约化字，所以群非交换。秩刻画了自由群的同构类型，却不能从“子群包含于一个低秩自由群”推出子群的秩也低；\secref{b1:sec:0804}将给出 \(F_2\) 的秩为三的子群。
\symidx{\(F_r\)：秩为 \(r\) 的自由群}

% !TeX root = ../../Book1.tex
% 纲要标题：### 8.2 群展示
% 纲要位置：计划纲要.md，第 341 行。

\section{群展示}
\label{b1:sec:0802}

自由群中的自由生成字母没有额外约束；有限循环群的生成元则满足有限阶关系，二面体群的旋转和反射还满足共轭关系。把这些等式写成关系字，再令关系字成为单位元，就能从自由群得到所需的商群。这里须特别区分两件事：具体元素满足所列关系，与这些关系已经完整描述该群，并非同一断言。后面的例子会分别核验这两步。

% 纲要内容（仅作写作计划，尚非教材正文）：
%
% * 生成元、关系与正规闭包
% * 从展示定义同态
% * 循环群、二面体群与四面体旋转群的展示
%

\subsection{生成元、关系与正规闭包}
\label{b1:subsec:0802:01}

施加关系时，设 \(K\) 为群，\(N\trianglelefteq K\)，\(q:K\to K/N\) 为商映射。若要求关系字 \(r\in K\) 满足 \(q(r)=1\)，则对任意 \(u\in K\) 都自动有
\[
q(uru^{-1})=q(u)q(r)q(u)^{-1}=1.
\]
它的逆以及这些共轭的有限乘积也都映为单位元。因此，令一个关系字成为单位元时，它的所有共轭也必然成为单位元；施加一组关系所需商去的是包含这些关系字的最小正规子群。

\begin{definition}\label{b1:def:normal-closure}
设 \(K\) 为群，\(R\subseteq K\)。令
\[
\NormalClosure{R}_{K}
\coloneqq
\bigcap_{\substack{N\trianglelefteq K\\R\subseteq N}}N.
\]
将这个包含 \(R\) 的最小正规子群称为 \(R\) 在环境群 \(K\) 中的\term[zhenggui-bibao]{正规闭包}[normal closure]。
\end{definition}
\symidx{\(\NormalClosure{R}_{K}\)：\(R\) 在 \(K\) 中的正规闭包}

正规闭包还有如下具体描述：
\[
\NormalClosure{R}_{K}
=\left\{\,
\prod_{i=1}^m u_i r_i^{\varepsilon_i}u_i^{-1}
\ \middle|\
\begin{gathered}
m\geq0,\quad u_i\in K,\quad r_i\in R,\\
\varepsilon_i\in\{1,-1\}\quad(1\leq i\leq m)
\end{gathered}
\,\right\}.
\]
当 \(m=0\) 时，所显示的乘积为空乘积，按 \(1_K\) 解释。两个这样的有限乘积相乘后仍属其中；取逆只须颠倒因子顺序并把各 \(\varepsilon_i\) 变号，所以右端是子群。用 \(v\in K\) 共轭时，只须把每个 \(u_i\) 换成 \(vu_i\)，因此它正规。这个子群包含 \(R\)，而任何包含 \(R\) 的正规子群都必须包含这些共轭因子及其有限乘积，所以它恰好等于定义中的交。

\ProseParagraphBreak
例如，在自由群 \(F=F\bigl(\{a,b\}\bigr)\) 中，普通生成子群 \(\langle b\rangle\) 只允许把 \(b\) 及其逆相乘，所得元素都是 \(b^k\)。正规闭包 \(\NormalClosure{b}_{F}\) 还允许先取任意共轭 \(ub^{\pm1}u^{-1}\)，再把这些共轭相乘。因此 \(aba^{-1}\) 属于 \(\NormalClosure{b}_{F}\)，却不属于 \(\langle b\rangle\)：它的约化字含有 \(a\) 字母，不可能等于任何 \(b^k\) 的约化字。所以 \(\langle b\rangle<\NormalClosure{b}_{F}\)。要使 \(b\) 在商群中成为单位元，也必须使它的所有共轭成为单位元。

\begin{definition}\label{b1:def:group-presentation}
设 \(S\) 为生成符号的集合，\(F(S)\) 为这些符号上的自由群，\(R\subseteq F(S)\) 为一组关系字。令
\[
\langle S\mid R\rangle
\coloneqq F(S)/\NormalClosure{R}_{F(S)}.
\]
若群 \(G\) 与这个商群同构，就称 \(\langle S\mid R\rangle\) 为 \(G\) 的一个\term[qun-zhanshi]{群展示}[group presentation]，也称群呈示；\(S,R\) 都有限时，称为有限展示。关系字 \(r\in R\) 表示所施加的等式 \(r=1\)。
\end{definition}
\symidx{\(\langle S\mid R\rangle\)：由生成元与关系给出的群展示}
\termidx[qun-chengshi]{群呈示}

本书约定把等式关系 \(v=w\) 写成关系字 \(vw^{-1}\)。在任意群中都有
\[
v=w\quad\Longleftrightarrow\quad vw^{-1}=1
\quad\Longleftrightarrow\quad v^{-1}w=1,
\]
所以用 \(v^{-1}w\) 也能表达同一个等式；所选次序只是统一记号的约定。

\ProseParagraphBreak
任何群都有展示。给定生成集 \(S\)，\cref{b1:thm:free-group-universal}把包含映射 \(S\to G\) 延拓为同态 \(\varphi:F(S)\to G\)；它的像包含生成集，故满射。取 \(R=\ker\varphi\)。因为核本身正规，所以 \(\NormalClosure{R}_{F(S)}=R\)，而\cref{b1:thm:first-isomorphism}给出 \(F(S)/R\cong\operatorname{im}\varphi=G\)，即 \(G\cong\langle S\mid R\rangle\)。这里的 \(R\) 把在 \(G\) 中取值为单位元的所有自由群元素都作为关系字纳入，只证明展示存在，并未解决怎样找出一组简便关系的问题。这既不保证生成集或关系集有限，也没有自动提供有效描述关系集或判别两个字是否相等的算法。

\subsection{展示与同态构造}
\label{b1:subsec:0802:02}

\begin{proposition}[展示的映射性质]\label{b1:prop:presentation-homomorphism}
设 \(S\) 为集合，\(R\subseteq F(S)\) 为一组关系字，并令 \(\langle S\mid R\rangle=F(S)/\NormalClosure{R}_{F(S)}\)。对任意群 \(H\)，从 \(\langle S\mid R\rangle\) 到 \(H\) 的群同态，与使所有关系字 \(r\in R\) 取值为 \(1_H\) 的集合映射 \(f:S\to H\) 一一对应。关系字的取值按各字母及其逆的像依次相乘得到。
\end{proposition}
\begin{proof}
由\cref{b1:thm:free-group-universal}，\(f\) 唯一延拓为同态 \(\widetilde f:F(S)\to H\)。记 \(N=\NormalClosure{R}_{F(S)}\)。满足全部关系的意思是对每个 \(r\in R\) 都有 \(\widetilde f(r)=1_H\)，即 \(R\subseteq\ker\widetilde f\)。核正规，故这又等价于 \(N\subseteq\ker\widetilde f\)。\cref{b1:thm:quotient-universal}因而给出唯一同态 \(\bar f:F(S)/N\to H\)，其公式为 \(\bar f(wN)=\widetilde f(w)\)。

反过来，若已给同态 \(\bar f\)，将它与商映射复合，再限制到 \(S\)，所得映射必使每个 \(r\in R\) 取值为目标群 \(H\) 的单位元。自由延拓与商映射分解的唯一性说明这两个过程互逆。
\end{proof}

验证一个群展示时，应当区分三个任务：检验所选元素满足关系，得到展示群 \(P\) 到目标群 \(G\) 的同态；进一步检验这些元素生成 \(G\)，才得到满同态，从而说明 \(G\) 是 \(P\) 的商群；最后证明这个满同态的核平凡，或直接构造逆同态，才能得到同构。概括地说，
\[
\begin{aligned}
\text{关系成立}&\Longrightarrow\text{指定取值延拓为同态};\\
\text{生成元的像生成目标}&\Longrightarrow\text{该同态满射};\\
\text{该同态满射且核平凡}&\Longrightarrow\text{该同态为同构}.
\end{aligned}
\]

核平凡不能单独保证同构：加法群同态 \(\mathbb Z\to\mathbb Z\)，\(n\mapsto2n\)，核平凡，却不是到所给目标群的同构。

\ProseParagraphBreak
关系成立且像生成目标群，也仍然不能保证同构。令 \(P=\langle a\mid a^6=1\rangle\)，并取已知的六阶循环群 \(C_6=\langle c\rangle\)。关系把 \(P\) 中每个字化为 \(1,a,\ldots,a^5\) 之一，所以 \(P\) 至多有六个元素；而 \(a\mapsto c\) 因 \(c^6=1\) 给出满同态 \(P\to C_6\)，所以这六个候选元素彼此不同，\(P\cong C_6\)。若改把 \(a\) 送到三阶循环群 \(C_3\) 的生成元 \(d\)，关系 \(d^6=1\) 仍成立，像也生成 \(C_3\)，却得到一个有非平凡核的满同态：\(a^3\ne1\)，而 \(d^3=1\)。

\ProseParagraphBreak
例如 \(P=\langle a,b\mid ab=ba\rangle\) 中，每个字可整理为 \(a^mb^n\)。映射 \(a\mapsto(1,0),b\mapsto(0,1)\) 给出 \(P\to\mathbb Z^2\)；反向映射 \((m,n)\mapsto a^mb^n\) 因交换关系而保持加法。两者复合都是恒等，故 \(P\cong\mathbb Z^2\)，且所列标准形唯一。若再加入 \(a^2=b^3=1\)，则得到 \(C_2\times C_3\)，而不是一个含有任意多交错字的非交换群。

\subsection{循环群、二面体群与四面体旋转群的展示}
\label{b1:subsec:0802:03}

有限循环群最简单的展示为
\[
C_n\cong\langle a\mid a^n=1\rangle\qquad(n\ge1).
\]
关系使每个字成为 \(1,a,\ldots,a^{n-1}\) 之一；映到已知的 \(n\) 阶循环群，说明这 \(n\) 个元素彼此不同。无限循环群展示则是一个生成元而没有关系。

\ProseParagraphBreak
这给出验证有限群展示的一种具体方法：先用关系把每个元素化为有限个候选标准形，若候选形式有 \(M\) 个，便得到展示群阶数至多为 \(M\) 的上界；再把展示群满射到一个已知 \(M\) 阶群。满射使展示群至少有 \(M\) 个元素，所以它恰有 \(M\) 个元素，该同态为同构，候选标准形也两两不同。下面的二面体群和四面体旋转群都用这种方法验证。

\begin{proposition}[二面体群的展示]\label{b1:prop:dihedral-presentation}
设 \(n\geq3\) 为整数。则
\[
D_{2n}\cong\langle r,s\mid r^n=1,\ s^2=1,\ srs^{-1}=r^{-1}\rangle.
\]
\end{proposition}
\begin{proof}
\cref{b1:prop:dihedral-normal-form}说明，具体旋转 \(r\) 与反射 \(s\) 满足这三条关系，生成 \(D_{2n}\)，而且 \(|D_{2n}|=2n\)。因此\cref{b1:prop:presentation-homomorphism}给出展示群到 \(D_{2n}\) 的同态，将两个生成符号送到相应对称；其像含 \(r,s\)，故为满射。

展示中的关系给出 \(sr^k=r^{-k}s\)，可把每个 \(s\) 推到字的最右端，再约去 \(s^2\) 及 \(r^n\)。于是每个元素都具有形式 \(r^i s^\varepsilon\)，其中 \(0\le i<n\)、\(\varepsilon\in\{0,1\}\)，展示群至多有 \(2n\) 个元素。满像已有 \(2n\) 个元素，故同态必为同构，标准形也唯一。
\end{proof}

正四面体旋转群的展示还可以从第六章的内半直积来理解。在 \(A_4\) 中，\(V_4=\{1,(12)(34),(14)(23),(13)(24)\}\) 是正规子群，因为共轭保持双换位。取 \(t_0=(123)\) 和 \(H_0=\langle t_0\rangle\)。\(V_4\) 的非单位元都有二阶，而 \(H_0\) 的非单位元都有三阶，所以 \(V_4\cap H_0=\{1\}\)，从而 \(V_4H_0\) 含有 \(4\cdot3=12\) 个元素，等于 \(A_4\)。由\cref{b1:thm:internal-semidirect}，\(A_4\cong V_4\rtimes C_3\)，作用由 \(t_0\) 的共轭给出。下面分别用 \(x,y\) 描述 \(V_4\)，用 \(t\) 描述三阶因子：\(x^2=y^2=1\)、\(xy=yx\) 规定四元因子的乘法，\(t^3=1\) 规定三阶因子，最后两条共轭关系规定两个因子如何相互作用。

\begin{example}[正四面体旋转群的展示]\label{b1:eg:tetrahedral-presentation}
正四面体的旋转群 \(A_4\) 有展示
\[
A_4\cong\left\langle x,y,t\ \middle|\
\begin{gathered}
x^2=y^2=t^3=1,\quad xy=yx,\\
txt^{-1}=y,\quad tyt^{-1}=xy
\end{gathered}\right\rangle.
\]
\end{example}
\begin{proof}[验证]
令 \(P\) 表示例中右侧的抽象展示群。先在具体群 \(A_4\) 中取
\[
x_0=(12)(34),\qquad y_0=(14)(23),\qquad t_0=(123).
\]
这里 \(x_0y_0=(13)(24)\)，所以 \(x_0,y_0\) 都有二阶且彼此交换，
\(V_0=\{1,x_0,y_0,x_0y_0\}\) 是四元子群。共轭只须把轮换中的各符号施以 \(t_0\)：得到
\[
t_0x_0t_0^{-1}=(23)(14)=y_0,
\qquad
t_0y_0t_0^{-1}=(24)(13)=x_0y_0.
\]
又因 \(t_0^3=1\)，这些置换满足全部关系。

三个集合 \(V_0,V_0t_0,V_0t_0^2\) 两两不交：若 \(vt_0^i=wt_0^j\)，其中 \(v,w\in V_0\)、\(0\le i,j<3\)，则 \(t_0^{i-j}\in V_0\)，但 \(V_0\) 没有三阶元素，所以 \(i=j\) 且 \(v=w\)。它们共含十二个元素，因而 \(x_0,y_0,t_0\) 生成整个 \(A_4\)。由\cref{b1:prop:presentation-homomorphism}，相应生成符号的取值给出满同态
\[
\varphi:P\longrightarrow A_4,
\qquad x\longmapsto x_0,\quad y\longmapsto y_0,\quad t\longmapsto t_0.
\]

现在在 \(P\) 中令 \(V=\langle x,y\rangle\)。由 \(x^2=y^2=1\) 和 \(xy=yx\)，\(V\) 的每个元素都可写成 \(x^\varepsilon y^\delta\)，其中 \(\varepsilon,\delta\in\{0,1\}\)，故 \(|V|\leq4\)。其余关系给出
\[
t x t^{-1}=y,\qquad
t y t^{-1}=xy,\qquad
t(xy)t^{-1}=y(xy)=x.
\]
这表明 \(t\) 的共轭把 \(x,y,xy\) 依次送到 \(y,xy,x\)，因而 \(tVt^{-1}=V\)。由于 \(t^{-1}=t^2\)，逆共轭也保持 \(V\)，具体有 \(t^{-1}xt=xy\)、\(t^{-1}yt=x\)。\(x,y\) 及其逆都属于 \(V\)，所以它们的共轭也保持 \(V\)。一般地，若子群在一组群生成元及其逆的共轭下保持，那么任意字的共轭也保持它，因为用乘积 \(uv\) 共轭就是先用 \(v\)、再用 \(u\) 共轭。因此 \(V\trianglelefteq P\)。商群 \(P/V\) 由 \(tV\) 生成，且 \((tV)^3=V\)，故 \([P:V]\leq3\)。于是
\[
|P|=|V|[P:V]\leq4\cdot3=12.
\]
此外，\(P/V\) 的三个候选元素为 \(V,tV,t^2V\)，而 \(V\) 正规，所以每个 \(P\) 中的元素都可写成
\[
x^\varepsilon y^\delta t^j,
\qquad \varepsilon,\delta\in\{0,1\},\quad 0\leq j<3.
\]
满同态 \(\varphi\) 的目标 \(A_4\) 有十二个元素，故上述上界迫使 \(|P|=12\)，并使 \(\varphi\) 成为同构。

几何上，绕顶点与对面中心的三分之一周转动给出三轮换，绕一对对边中点连线的半周转动给出双换位，故这些确实是四面体的旋转。
\end{proof}

同一个群可以有不同的展示。\cref{b1:eg:tetrahedral-presentation}保留 \(V_4\) 的两个生成元，使半直积结构清楚可见；为了追求最少的生成元，可以用 \(y=txt^{-1}\) 消去 \(y\)，但关系会变长。展示的简短程度、标准形的易算程度及结构的可见程度，往往需要分别考虑。

% !TeX root = ../../Book1.tex
% 纲要标题：### 8.3 自由积
% 纲要位置：计划纲要.md，第 347 行。

\section{自由积}
\label{b1:sec:0803}

把两个群组成一个新群，可以提出不同要求。直积不仅保留两个因子的乘法，还让来自不同因子的元素交换；自由积只保留各因子内部已有的乘法。本节仍用字来构造群，不过现在同一因子中的相邻字母要先相乘，不同因子的字母则保持次序。约化字将解释自由积的泛性质，也说明只要两个因子都非平凡，它们的自由积便是无限群。

% 纲要内容（仅作写作计划，尚非教材正文）：
%
% * 自由积的约化字
% * 自由积的泛性质
% * 直积与自由积的区别
%

\subsection{自由积的约化字}
\label{b1:subsec:0803:01}

自由群的字母来自所选生成符号及其形式逆；自由积的字母则取自两个因子中的全部非单位元素，一个字母本身就可能是在因子内由许多生成元相乘得到的元素。即使两个群的底集合有重叠，也先给元素加上所属因子的标记。这样，同一名称来自不同因子时仍是不同字母，字的化简才不会误用另一因子的乘法。

\begin{definition}\label{b1:def:free-product-word}
设 \(G,H\) 为群。分别取 \(G\setminus\{1_G\}\) 与 \(H\setminus\{1_H\}\) 的不相交副本作为字母表。由这些字母组成的有限序列称为自由积字；若每对相邻字母都来自不同因子，则称为约化字。长度为零的空字也约化。
\end{definition}

从所有有限字出发，允许在同一因子内把相邻的 \(x,y\) 换为它们的乘积 \(xy\)，乘积为单位元时删去这两个字母；也允许这些操作的逆操作。由此生成的等价关系与拼接相容。与自由群约化相比，自由群只消去相邻互逆字母，自由积则把相邻同因子字母先在该因子内相乘；只有乘积为单位元时才消去。因此，下面仍从左到右保存一个约化字，但读取同因子的新字母时，须用因子内的乘法更新末尾。

\begin{lemma}\label{b1:lem:free-product-normal-form}
设 \(G,H\) 为群。在以 \((G\setminus\{1_G\})\sqcup(H\setminus\{1_H\})\) 为字母表的全部有限字上，令 \(\sim\) 为由合并或拆分同一因子中的相邻字母所生成的等价关系，其中乘积为单位元时删除相应字母。则每个等价类中恰有一个约化字。
\end{lemma}
\begin{proof}
从左到右读取字母，始终保存一个约化字。若当前字为空，就接入新字母；若当前字不为空且新字母与末尾来自不同因子，就接在末尾；若来自同一因子，就在该因子内相乘，乘积非单位元时用乘积替换末尾，乘积为单位元时删去末尾。这给出一个与原字等价的约化字。

验证输出不随允许的操作改变。设 \(x,y\) 来自同一因子 \(G\)，比较连续读取 \(x,y\) 与一次读取 \(xy\)；当 \(xy=1_G\) 时，后一操作指不作任何改变。若原字为空或末尾来自 \(H\)，两次读取先添 \(x\) 再合并成 \(xy\)，与直接读取相同。

若原末尾为 \(z\in G\)，且 \(zx\ne1\)，两次读取按 \((zx)y=z(xy)\) 合并；乘积为单位元时删去末尾，否则保留乘积。若 \(xy=1\)，这就恢复 \(z\)。若 \(zx=1\)，第一次删去 \(z\)，前一字母如存在必来自 \(H\)，第二次便接上 \(y\)；直接读取的结果由 \(z(xy)=y\) 给出，也一样。交换 \(G,H\) 得到另一因子的同样结论。

因此一次合并或拆分同因子字母不改变上述约化过程的输出。约化字的输出是自身，所以同一等价类不能包含两个不同约化字。
\end{proof}

设 \(G,H\) 为群，并在\cref{b1:def:free-product-word}的自由积字上取由同因子相邻字母合并与拆分生成的等价关系。等价类在拼接下组成群：结合律从拼接继承，空字为单位，逆由字母倒序取逆得到。将此群称为 \(G,H\) 的\term[ziyou-ji-product]{自由积}[free product]，记为 \(G*H\)。
\symidx{\(G*H\)：两个群的自由积}

\ProseParagraphBreak
记空字为 \(\varepsilon\)，字 \(w\) 的等价类为 \([w]\)。如果把因子的单位元也保留为长度一的约化字，它就会与空字重复表示自由积的单位元。因此字母表一开始就排除了 \(1_G,1_H\)，两个因子的嵌入必须把单位元送到空字类，而非单位元才对应长度为一的字。具体地，定义
\[
j_G:G\longrightarrow G*H,\qquad
j_G(g)=
\begin{cases}
[\varepsilon],&g=1_G,\\
[g],&g\ne1_G.
\end{cases}
\]
同样定义 \(j_H:H\rightarrow G*H\)，其中 \(j_H(1_H)=[\varepsilon]\)，非单位元 \(h\) 的像为 \([h]\)。

\ProseParagraphBreak
若 \(x,y\in G\) 中有一个为单位元，拼接空字不改变字，故 \(j_G(xy)=j_G(x)j_G(y)\)。若两者都非单位元，同因子字母的合并规则给出同一等式；尤其 \(xy=1_G\) 时，字 \(xy\) 删除为空字，所以也满足同态条件。不同非单位元给出不同的长度为一的约化字，且都不同于空字；由\cref{b1:lem:free-product-normal-form}，\(j_G\) 单射。对 \(H\) 同样成立。因两种字母保留所属因子的标记，两个像的交只有单位元。这些论证也包括某个因子为平凡群的情形。

\ProseParagraphBreak
约化标准形在集合论上给出了比较元素的方法：两个元素相等，当且仅当它们的约化字逐字相同。不过，要把上述从左到右的约化过程实现为处理有限输入的有效算法，还须能有效表示和比较各因子的元素、计算同一因子内的乘积，并判定所得乘积是否为单位元。由 \(j_G\) 的单射性，只使用 \(G\) 中元素写成的字在 \(G*H\) 中表示单位元，当且仅当它在 \(G\) 中的乘积为单位元。因此，因子的字问题也包含在自由积的字问题中，不能仅凭约化标准形的唯一性，就断言任意自由积的字问题可解。

\ProseParagraphBreak
例如在 \(C_2*C_3\) 中，令 \(a,b\) 分别为两个循环因子的生成元，\(a\) 的阶为 \(2\)，\(b\) 的阶为 \(3\)。这里要区分按因子元素计数的自由积字长度（也叫音节长度），与继续展开成所选生成符号后所得字的字母个数。\(b^2\) 是一个因子中的非单位元素，因而在自由积中是一个字母：\(ab^2a\) 是长度为 \(3\) 的自由积约化字，而所写展开字 \(abba\) 有 \(4\) 个字母。后一个数取决于所写表达，不能直接当作元素相对于生成元及其逆的最短字长；例如 \(b^2=b^{-1}\)，同一元素也可写成 \(ab^{-1}a\)。由 \(a^2=b^3=1\) 得
\[
(ab^2a)(ab)=ab^2b=a,
\]
但 \((ab)^k\) 对每个 \(k\ge1\) 都是长度 \(2k\) 的约化字，所以 \(ab\) 具有无限阶。更一般地，只要 \(G,H\) 都非平凡，任取 \(1_G\neq g\in G\)、\(1_H\neq h\in H\)，则每个 \((gh)^k\) 都是长度为 \(2k\) 的约化字。由\cref{b1:lem:free-product-normal-form}，不同长度的这些约化字表示不同元素，故 \(G*H\) 为无限群。特别地，两个非平凡有限群的自由积仍然无限；若某个因子平凡，自由积则同构于另一个因子。

\subsection{自由积的泛性质}
\label{b1:subsec:0803:02}

\begin{theorem}[自由积的泛性质]\label{b1:thm:free-product-universal}
设 \(G,H,K\) 为群，\(G*H\) 为 \(G,H\) 的自由积，\(j_G:G\to G*H\)、\(j_H:H\to G*H\) 为自然嵌入。给定群同态 \(f:G\to K\)、\(g:H\to K\)，存在唯一群同态 \(u:G*H\to K\)，使 \(uj_G=f\)、\(uj_H=g\)，即在两个因子上的限制分别为 \(f,g\)。
\end{theorem}
\begin{proof}
把一个字映为各字母像的乘积，并把空字映为 \(1_K\)。同因子字母的合并不改变这个乘积，因为 \(f,g\) 各自保持乘法；当乘积为因子的单位元时，其像为 \(1_K\)，删去这两个字母也不改变取值。所以映射在等价类上良定义，并保持拼接乘法，且满足 \(uj_G=f\)、\(uj_H=g\)。两个因子共同生成 \(G*H\)，故这样的同态唯一。
\end{proof}

自由积的泛性质没有要求 \(f(G)\) 与 \(g(H)\) 交换。例如从 \(C_2*C_3\) 到 \(S_3\)，可以指定 \(a\mapsto(12)\)、\(b\mapsto(123)\)，得到满同态。若再令 \((ab)^2=1\)，商群就同构于 \(S_3\)：该关系给出 \(aba=b^{-1}\)，即 \(ab=b^{-1}a\)，所以每个字化为 \(b^i a^\varepsilon\)，其中 \(i\in\{0,1,2\}\)、\(\varepsilon\in\{0,1\}\)，至多六个元素；相应置换像已经有六个元素。

\ProseParagraphBreak
若群 \(P\) 及同态 \(i_G:G\to P\)、\(i_H:H\to P\) 也满足\cref{b1:thm:free-product-universal}的泛性质，则存在唯一同构 \(G*H\to P\)，使它与 \(j_G,j_H\) 的复合分别为 \(i_G,i_H\)。确实，两个泛性质分别给出两个方向的延拓；两次复合都固定相应因子的像，故由延拓唯一性为恒等。这与\cref{b1:prop:free-universal-unique}的双向映射论证相同。

\ProseParagraphBreak
给定任意群 \(K\)，从 \(S,T\) 到 \(K\) 的集合映射可以彼此独立地选择，再分别唯一延拓为 \(F(S),F(T)\) 到 \(K\) 的同态。自由积的泛性质把这两个同态唯一延拓为 \(F(S)*F(T)\to K\)。因此 \(F(S)*F(T)\) 与 \(F(S\sqcup T)\) 满足相同的自由群泛性质，特别地，有
\[
F(S\sqcup T)\cong F(S)*F(T)
\]
（这里取不相交副本）。具体地，把 \(S,T\) 的字母送入自由积中对应的因子，\cref{b1:thm:free-group-universal}给出 \(u:F(S\sqcup T)\to F(S)*F(T)\)。反向先把 \(S,T\) 分别送到 \(F(S\sqcup T)\) 的同名字母，用同一定理得到两个因子的同态，再由\cref{b1:thm:free-product-universal}得到 \(v:F(S)*F(T)\to F(S\sqcup T)\)。\(vu\) 固定所有字母，\(uv\) 固定两个因子的全部字母，所以两个复合都固定生成集，均为恒等。

\ProseParagraphBreak
若 \(G=\langle S\mid R\rangle\)、\(H=\langle T\mid U\rangle\)，且生成符号互不相交，则
\begin{equation}\label{b1:eq:free-product-presentation}
G*H\cong\langle S\sqcup T\mid R\cup U\rangle.
\end{equation}
记右侧为 \(P\)。其生成元在 \(P\) 中分别满足 \(R,U\)，由\cref{b1:prop:presentation-homomorphism}得到 \(G\to P,H\to P\)，再由\cref{b1:thm:free-product-universal}得到 \(u:G*H\to P\)。另一方面，两个因子在 \(G*H\) 中的生成元满足 \(R\cup U\)，展示的同态构造又给出 \(v:P\to G*H\)。两个复合在 \(S\sqcup T\) 的像上为恒等，故互为逆。这个描述没有额外加入交换关系 \([s,t]=sts^{-1}t^{-1}=1\)；下面证明加入跨因子交换关系以后得到直积。

\subsection{直积与自由积的区别}
\label{b1:subsec:0803:03}

自由积只保留两个因子各自内部的关系；直积比它多出的全部关系，正是来自两个因子的元素必须彼此交换。下面的核公式说明，商去表示这些交换关系的关系字所生成的正规闭包，就得到直积。

\begin{proposition}\label{b1:prop:product-versus-free-product}
设 \(G,H\) 为群，\(G*H\) 为其自由积，并通过自然嵌入把 \(G,H\) 视为 \(G*H\) 的子群。自然满同态
\[
\pi:G*H\longrightarrow G\times H,\qquad
g\longmapsto(g,1),\quad h\longmapsto(1,h)
\]
满足
\[
\ker\pi
=\NormalClosure{\bigl\{\,[g,h]\mid g\in G,\ h\in H\,\bigr\}}_{G*H},
\]
其中 \([g,h]=ghg^{-1}h^{-1}\)，右端是\cref{b1:def:normal-closure}定义的、在环境群 \(G*H\) 中取得的正规闭包。
\end{proposition}
\begin{proof}
由\cref{b1:thm:free-product-universal}，两个因子嵌入直积的同态延拓为 \(\pi\)；每个 \((g,h)=(g,1)(1,h)\) 在其像中，所以 \(\pi\) 满射。

记 \(N=\NormalClosure{\bigl\{\,[g,h]\mid g\in G,\ h\in H\,\bigr\}}_{G*H}\)，商映射为 \(q:G*H\to Q=(G*H)/N\)。每个 \([g,h]\) 都被 \(\pi\) 送到单位，而 \(\ker\pi\) 正规，所以正规闭包的定义给出 \(N\subseteq\ker\pi\)。由\cref{b1:thm:quotient-universal}，存在唯一同态 \(\bar\pi:Q\to G\times H\) 满足 \(\bar\pi q=\pi\)。

在 \(Q\) 中两个因子的像逐元素交换，因此
\[
G\times H\longrightarrow Q,\qquad
(g,h)\longmapsto \overline g\,\overline h
\]
给出同态 \(\sigma\)：确实，\((\overline g\,\overline h)(\overline {g'}\,\overline {h'})=\overline {gg'}\,\overline {hh'}\)，因为中间的 \(\overline h,\overline {g'}\) 可交换。

于是 \(\bar\pi\sigma(g,h)=(g,h)\)，而 \(\sigma\bar\pi\) 固定 \(Q\) 中两个因子的像；这些像生成 \(Q\)，故 \(\sigma\bar\pi=\operatorname{id}_Q\)。所以 \(\bar\pi\) 是同构，\(\ker\pi=\ker(\bar\pi q)=\ker q=N\)。
\end{proof}

若 \(S,T\) 分别生成 \(G,H\)，则只需施加生成元之间的交换关系。令
\[
N_0=\NormalClosure{\bigl\{\,[s,t]\mid s\in S,\ t\in T\,\bigr\}}_{G*H}.
\]
显然 \(N_0\subseteq\ker\pi\)。在 \((G*H)/N_0\) 中，\(S\) 与 \(T\) 的像逐元素交换；它们的逆元也彼此交换，逐项换序便知 \(G\) 中任意元素的像都与 \(H\) 中任意元素的像交换。因此，\cref{b1:prop:product-versus-free-product}中所有交换子 \([g,h]\) 都属于 \(N_0\)，从而 \(\ker\pi\subseteq N_0\)，故 \(N_0=\ker\pi\)。在\eqref{b1:eq:free-product-presentation}的展示中，将这些生成元交换子作为新关系字，正对应于商去 \(N_0\)。因而若 \(G=\langle S\mid R\rangle\)、\(H=\langle T\mid U\rangle\)，且两组生成符号互不相交，就得到
\[
G\times H\cong
\left\langle S\sqcup T\ \middle|\ R\cup U\cup
\{\,[s,t]\mid s\in S,\ t\in T\,\}\right\rangle.
\]

当两个因子都非平凡时，选非单位元 \(g\in G,h\in H\)，字 \(ghg^{-1}h^{-1}\) 长四且约化，由\cref{b1:lem:free-product-normal-form}，它不是单位，却被 \(\pi\) 送到单位，故自然同态不单射。直积中的相应交换子却必为单位。

\ProseParagraphBreak
特别地，\(C_2*C_2\) 是无限二面体群。设两个二阶生成元为 \(a,b\)，令 \(r=ab,s=a\)。则
\[
srs^{-1}=r^{-1},
\]
而 \(r\) 具有无限阶。每个偶数长度的交错字为 \((ab)^k\) 或 \((ba)^k=r^{-k}\)（\(k\ge0\)）。若 \(w\) 是奇数长度的约化字，在右端再乘 \(a=s\)，再约化后长度为偶数：末字母为 \(a\) 时消去末尾的 \(aa\)，末字母为 \(b\) 时则添上一个 \(a\)。因此 \(wa=r^k\) 对某个 \(k\in\mathbb Z\) 成立，再由 \(a^2=1\) 还原得
\[
w=(wa)a=r^ka=r^ks.
\]
这些形式两两不同：\(r^k\) 的约化字长度为偶数，\(r^ks\) 的约化字长度为奇数；同一种形式的不同指数若相等，则给出 \(r\) 的非零幂为一，与\cref{b1:lem:free-product-normal-form}矛盾。因此每个元素唯一写成 \(r^k\) 或 \(r^ks\)，其中 \(k\in\mathbb Z\)。它可具体作用在整数直线上：令 \(a\) 作用为 \(m\mapsto-m\)，\(b\) 作用为 \(m\mapsto1-m\)，则 \(r\) 作用为平移 \(m\mapsto m-1\)。这同时展示了反射的几何来源与约化字的代数判别。对比 \(C_2\times C_2\)，后者只有四个元素，两次反射的乘积也只有二阶。

\ProseParagraphBreak
\cref{b1:tab:group-construction-maps}比较从四种构造对象映到目标群 \(K\) 的同态：可以指定哪些数据，以及这些数据应在 \(K\) 中满足什么条件。表中的直积一行讨论的是 \(G\times H\to K\)，其条件是两个因子的像逐元素交换；从 \(K\) 映到直积的同态则由两个坐标同态 \(K\to G\)、\(K\to H\) 决定。

\begin{table}[htbp]
\centering
\caption[四种群构造的同态条件]{从四种群构造到目标群的同态条件。目标 \(K\) 为任意群；满足对应条件时，给定数据唯一延拓为所列群到 \(K\) 的同态。}
\label{b1:tab:group-construction-maps}
\begin{tabularx}{\linewidth}{@{}>{\raggedright\arraybackslash}p{0.28\linewidth}>{\raggedright\arraybackslash}p{0.27\linewidth}>{\raggedright\arraybackslash}X@{}}
\toprule
映到 \(K\) 的源群及其映射性质 & 给定的数据 & 目标中的相容条件 \\
\midrule
自由群 \(F(S)\)\newline
\cref{b1:thm:free-group-universal}
& 集合映射 \(S\to K\)
& 无额外关系 \\
\addlinespace
展示群 \(\langle S\mid R\rangle\)\newline
\cref{b1:prop:presentation-homomorphism}
& 集合映射 \(S\to K\)
& 每个关系字在 \(K\) 中取值为单位元 \\
\addlinespace
自由积 \(G*H\)\newline
\cref{b1:thm:free-product-universal}
& 同态 \(G\to K\)、\(H\to K\)
& 两个像之间无额外相容条件 \\
\addlinespace
直积 \(G\times H\)\newline
\cref{b1:prop:finite-product-maps}
& 同态 \(G\to K\)、\(H\to K\)
& 两个像逐元素交换 \\
\bottomrule
\end{tabularx}
\end{table}

% !TeX root = ../../Book1.tex
% 纲要标题：### 8.4* 自由群的子群与图方法
% 纲要位置：计划纲要.md，第 353 行。

\OptionalSection{自由群的子群与图方法}{b1:sec:0804}

自由群中的子群未必能靠删去几个原有基元素得到；它甚至可能需要比母群更多的自由生成元。本节把子群的字解释为陪集图中的闭路径，用一棵生成树选出新的自由基。所需的图论事实会在证明中说明，不预设拓扑空间或基本群的知识。有限指数时，基的个数可由边数与顶点数算出；节末再讨论这些构造能够解决哪些字的判定问题。

% 纲要内容（仅作写作计划，尚非教材正文）：
%
% * 图的闭路径群与生成树
% * Schreier 图与子群自由性
% * 有限指数子群与 Schreier 秩公式
% * 字问题与 Cayley 图初步
%
% ---
%

\subsection{图的闭路径群与生成树}
\label{b1:subsec:0804:01}

为研究自由群的子群，先处理一个组合问题：在连通图中，怎样用一族选定的闭路径表达所有闭路径，并辨认其中的关系？

\ProseParagraphBreak
本节的图 \(\Gamma\) 由顶点和连接顶点的几何边组成，并约定至少有一个顶点。允许两个顶点之间有多条边，也允许边的两端是同一个顶点，这样的边称为环。每条几何边有两个相反方向 \(e,\bar e\)，即使是环也区分方向，并规定 \(\overline{\bar e}=e\)。子图选取部分顶点和几何边，且每条所选边的两个端点都须保留。

\ProseParagraphBreak
一条边路径是有限个首尾相接的有向边组成的序列；在每个顶点还允许不经过任何边的空路径。路径 \(e_1\cdots e_k\) 的反向路径为 \(\bar e_k\cdots\bar e_1\)，记作 \((e_1\cdots e_k)^{-1}\)。不含相邻反向边 \(e\bar e\) 的路径称为约化路径。插入或删除这样的一对边称为一次往返变换，由有限次往返变换得到的路径视为等价；它们有相同的起点和终点。

\ProseParagraphBreak
若任意两顶点间都有边路径，则图称为连通的。正长度的约化闭路径若除首尾外不重复顶点，称为圈。连通而没有圈的图称为树；包含 \(\Gamma\) 全部顶点的一棵树子图称为 \(\Gamma\) 的生成树。树中的约化路径不会重复顶点，否则取两个相同顶点之间最短的一段就会得到圈。任意边路径则仍可含往返。例如在只有一条边 \(e:o\rightarrow v\) 的树中，\(e\) 与 \(e\bar e e\) 是不同的路径序列，但往返等价；只有前者约化。

\ProseParagraphBreak
设 \(\Gamma\) 为图，\(o\) 为其一个顶点。从 \(o\) 回到 \(o\) 的闭路径在往返变换的等价关系下，以拼接为乘法组成群；将这个群称为图 \(\Gamma\) 在基点 \(o\) 处的\term[tu-de-jibenqun]{基本群}[fundamental group]，记为 \(\pi_1(\Gamma,o)\)。这里采用边路径的组合定义，不预设一般拓扑空间的基本群理论。拼接与往返变换相容，因而在等价类上良定义；结合律由拼接给出，\(o\) 处的空路径为单位，逆由反向路径给出。
\symidx{\(\pi_1(\Gamma,o)\)：图的基本群}

\ProseParagraphBreak
为了比较散布在不同顶点的边，需要先把它们都接到同一个基点。生成树恰好提供从基点到每个顶点的唯一约化路径，而且树内闭路径都能往返消去。借助这些路径，每条非树边都能补成基点处的闭路径；树只负责连接顶点，余下的边才记录闭路径的自由选择。

\cref{b1:fig:graph-tree-loop}中，从 \(o\) 到 \(v\) 和从 \(o\) 到 \(w\) 的两条树路径共用 \(o\) 到 \(u\) 的一段；构造并不要求它们只在基点相交。

\begin{figure}[htbp]
\centering
\begin{tikzpicture}[scale=0.9,>=stealth]
\node[circle,draw,inner sep=2pt] (o) at (0,0) {\(o\)};
\node[circle,draw,inner sep=2pt] (u) at (2,0) {\(u\)};
\node[circle,draw,inner sep=2pt] (v) at (4,1.25) {\(v\)};
\node[circle,draw,inner sep=2pt] (w) at (4,-1.25) {\(w\)};
\draw[very thick] (o)--(u)--(v);
\draw[very thick] (u)--(w);
\draw[->,thick,dashed,color=AlgebraBlue] (v) to[bend left=30] node[right] {\(e\)} (w);
\node[above] at (3,0.63) {\(p_v\)};
\node[below] at (3,-0.63) {\(p_w^{-1}\)};
\end{tikzpicture}
\caption[生成树与非树边构成的基环路]{实线是生成树中的边，虚线是非树边 \(e:v\to w\)。树路径 \(p_v:o\to u\to v\) 与反向树路径 \(p_w^{-1}:w\to u\to o\) 共用 \(o,u\) 之间的边；对应的自由基元素是闭路径类 \([p_v e p_w^{-1}]\)。}
\label{b1:fig:graph-tree-loop}
\end{figure}

\begin{lemma}[生成树给出的闭路径自由基]\label{b1:lem:graph-loop-free}
设 \(\Gamma\) 是连通图，\(T\subseteq\Gamma\) 是一棵生成树，\(o\) 是指定基点。为每条不在 \(T\) 中的几何边选一个方向，组成集合 \(E_0\)。从 \(o\) 到顶点 \(v\) 的唯一约化树路径记为 \(p_v\)，其中 \(p_o\) 为空路径。则 \(\pi_1(\Gamma,o)\) 是自由群，自由基由下列闭路径的等价类组成：
\[
\ell_e=[p_v\,e\,p_w^{-1}]\qquad
(e\in E_0,\ e:v\longrightarrow w).
\]
\end{lemma}

\begin{proof}
树的两顶点间有唯一约化路径。连通性先给出一条有限边路径，反复删除相邻往返就得到约化路径，所以存在。若有两条不同的约化路径，因它们都不重复顶点，取首次分离与再次汇合的区段便得到圈，违背树的无圈性。特别地，树内的任何闭路径均可消为空路径；否则最终留下的非空约化闭路径会包含圈。

把 \(E_0\) 当作自由字母，由\cref{b1:thm:free-group-universal}把集合映射 \(e\mapsto\ell_e\) 延拓为同态 \(\Phi:F(E_0)\rightarrow\pi_1(\Gamma,o)\)。再定义读取映射 \(\Psi:\pi_1(\Gamma,o)\rightarrow F(E_0)\)：沿闭路径读取边，树边略过，非树边按所选方向读作 \(e\)，沿相反方向读作 \(e^{-1}\)，最后约化。插入往返时，若是树边不会读出字母，若是非树边则增加相邻逆字母，因此 \(\Psi\) 良定义；拼接路径对应拼接标签字，所以它是同态。

沿基环路 \(p_v e p_w^{-1}\) 应用读取映射，只得到字母 \(e\)，故 \(\Psi\Phi\) 在所有自由字母上为恒等，因而是恒等同态。为检查另一方向，设闭路径依次走边 \(e_i:v_{i-1}\rightarrow v_i\)，其中 \(v_0=v_k=o\)。在每个相邻位置插入 \(p_{v_i}^{-1}p_{v_i}\)，得到它等价于
\[
\prod_{i=1}^k
\bigl(p_{v_{i-1}}e_i p_{v_i}^{-1}\bigr).
\]
树边对应的括号项是树内闭路径，可消为空路径；非树边对应基环路或其逆。因此读出字母再送回环路，恰好恢复原来的等价类，即 \(\Phi\Psi\) 也是恒等。两个同态互逆。
\end{proof}

连通图总可选一棵生成树。有限图可以从根开始，每次添加一条通向新顶点的边；只要还有未加入的顶点，连通性就保证存在这样的边。每次只接入一个新顶点，因此不会产生圈，有限次后便得到生成树。

\ProseParagraphBreak
一般无限图的存在性证明使用通常的选择公理。考虑包含全部顶点、没有圈的子图，按边集包含排序；只有顶点而没有边的子图说明这个偏序集非空。每条链的并仍没有圈：若并中有圈，圈只用有限条边，链中必有一个子图已包含全部这些边，矛盾。因此链的并是同一偏序集中的上界，\cref{b1:thm:A-zorn}给出一个极大的无圈子图。若它不连通，原图中连接两个连通分支的路径必有一条跨分支边；加进这条边仍不产生圈，却严格增大该子图，矛盾。所以极大子图连通，是生成树。只讨论有限图的例子时，可直接使用上一段的逐步构造。

\subsection{Schreier 图与子群自由性}
\label{b1:subsec:0804:schreier-subgroup-freeness}

现在把自由群的字解释为路径。对自由群的子群 \(H\)，图的顶点记录右陪集，字母记录从一个陪集到另一个陪集的移动；从顶点 \(H\) 出发并回到 \(H\) 的字恰好给出 \(H\) 中的元素。

\ProseParagraphBreak
\term[nielsen-schreier-dingli]{Nielsen--Schreier 定理}[Nielsen--Schreier theorem]断言自由群的每个子群仍是自由群。生成树先把 Schreier 图的闭路径群写成自由群，路径标签再把这个闭路径群识别为原来的子群。\footnote{尼尔森（Jakob Nielsen，1890-1959），丹麦数学家，研究自由群的子群和曲面的拓扑变换。}

\begin{theorem}[Nielsen--Schreier]\label{b1:thm:nielsen-schreier}
设 \(F\) 为自由群，\(H\leq F\)。则 \(H\) 也是自由群。
\end{theorem}

\begin{proof}
设 \(H\le F(S)\)。构造以右陪集 \(Hg\) 为顶点的\term[schreier-peiji-tu]{Schreier 陪集图}[Schreier coset graph] \(\Gamma_H\)：对每个顶点 \(Hg\) 和 \(s\in S\)，画一条指定正向为
\[
Hg\xrightarrow{s}Hgs
\]
的几何边，其反向标记为 \(s^{-1}\)。若 \(Hg=Hg_1\)，则 \(g_1=hg\) 对某个 \(h\in H\) 成立，因而 \(Hg_1s=Hhgs=Hgs\)，所以边的终点与代表元选择无关。这里不要求 \(H\) 正规，也没有把这些陪集当作商群作乘法。每条正向边由起点和基字母共同确定；即使两条边有同样的端点，也保留其不同身份。

对给定顶点 \(Hg\) 和 \(s\in S\)，唯一标记为 \(s^{-1}\) 的出边，是正向边
\[
Hgs^{-1}\xrightarrow{s}Hg
\]
的反向。因而对每个正字母或负字母，每个顶点都恰有一条相应出边。选择右陪集，是为了使从左至右读取字母与右乘相容：从 \(Hg\) 依次读取 \(a_1,\ldots,a_m\) 后，终点为 \(Hga_1\cdots a_m\)。特别地，任意字从 \(H\) 出发都确定唯一的带标签路径，任意顶点 \(Hg\) 都可这样到达，所以图连通。

以 \(H\) 为根。闭路径的标签乘积属于 \(H\)；往返变换只插入或删除互逆标签，不改变该乘积。因此标签乘积给出同态 \(\lambda_H:\pi_1(\Gamma_H,H)\rightarrow H\)。每个 \(h\in H\) 的字都有回到根的路径，故 \(\lambda_H\) 满射。

若闭路径的标签在自由群中为单位，\cref{b1:lem:unique-reduced-word}说明其唯一约化字为空；该引理的约化算法只删除相邻逆字母，因此可以逐次把这个标签字消为空字。图中每个顶点都有唯一的给定字母出边，所以路径上相邻逆字母必沿同一条几何边往返：第一条边的反向已经是从中间顶点出发、具有所需逆标签的唯一边。字的每次消去因此对应路径的一次往返消去，最终路径也消为空路径，故 \(\lambda_H\) 单射。于是得到标签同构
\begin{equation}\label{b1:eq:schreier-label-isomorphism}
\lambda_H:\pi_1(\Gamma_H,H)\xrightarrow{\sim}H.
\end{equation}

选一棵生成树，\cref{b1:lem:graph-loop-free}给出第一个同构，式\eqref{b1:eq:schreier-label-isomorphism}给出第二个同构：
\[
F(E_0)\xrightarrow[\text{生成树}]{\sim}
\pi_1(\Gamma_H,H)\xrightarrow[\text{标签}]{\lambda_H}H.
\]
前一个同构证明闭路径群自由，后一个同构把它识别为所研究的子群。若树路径 \(p_v\) 的标签字为 \(t_v\)，基元素可显式写为
\begin{equation}\label{b1:eq:schreier-basis}
t_v s\,t_w^{-1}.
\end{equation}
这里 \(v\xrightarrow{s}w\) 取遍不在树中的正向边。两条到达 \(w\) 的路径 \(t_vs,t_w\) 表示相同右陪集，所以
\[
Ht_vs=Ht_w\quad\Longrightarrow\quad t_vs\,t_w^{-1}\in H.
\]
这既验证了基元素属于 \(H\)，也确定了公式中两个树代表的次序。
\end{proof}

这个证明说明，一旦完整的 Schreier 图 \(\Gamma_H\) 和一棵生成树已经给出，就可由非树边读出自由基；它并不自动给出从任意子群描述计算全部陪集图的算法。事实上，识别两个顶点需要判定
\[
Hg=Hk\quad\Longleftrightarrow\quad gk^{-1}\in H,
\]
这涉及子群成员判定，并不等同于判断自由群元素 \(gk^{-1}\) 是否为单位。能够直接画出图的具体例子，另外使用了可计算的同态或陪集模型。

\ProseParagraphBreak
子群中的字关系在图中对应往返消去。选取生成树以后，树内闭路径本来就与空路径等价；利用这些等价关系，任意闭路径可以写成非树边所给基环路的乘积。这个过程只改变路径类的表达方式，没有给 \(H\) 添加新的关系。不同生成树可以给出不同的自由基，但标签同构\eqref{b1:eq:schreier-label-isomorphism}始终描述同一个子群。

\subsection{有限指数子群与 Schreier 秩公式}
\label{b1:subsec:0804:02}

\term[schreier-zhi-gongshi]{Schreier 秩公式}[Schreier rank formula]把有限指数子群的自由基个数化为图的边数与顶点数之差。

\begin{theorem}[Schreier 秩公式]\label{b1:thm:schreier-rank}
设 \(r\geq0\) 为整数，\(F_r\) 为秩为 \(r\) 的自由群，且子群 \(H\le F_r\) 的指数为有限数 \(d\)。则
\[
\operatorname{rank}H=1+d(r-1).
\]
\end{theorem}

\begin{proof}
Schreier 图有 \(d\) 个顶点，每个顶点对于每个基字母恰有一条正向边，故共有 \(dr\) 条几何边；反向边不重复计数。用前面的有限生成树构造选取 \(T\)：从一个根顶点开始，每一步同时添加一条边和一个尚未加入的新顶点；到达全部 \(d\) 个顶点时，恰好添加了 \(d-1\) 条边。因此 \(T\) 有 \(d-1\) 条几何边，非树边数为 \(dr-(d-1)\)。根据\cref{b1:thm:nielsen-schreier}证明中的构造，每条非树边由\eqref{b1:eq:schreier-basis}对应一个自由基元素；再由\cref{b1:prop:free-rank-invariant}，这一基的大小就是 \(H\) 的秩。\(r=0\) 时只能有 \(d=1\)，公式仍成立。
\end{proof}

例如 \(F_2=F\bigl(\{a,b\}\bigr)\) 到加法群 \(C_2\) 的同态 \(a\mapsto1,b\mapsto0\) 是满射，核 \(H\) 的指数为二。\cref{b1:fig:schreier-index-two}有两个顶点 \(H,Ha\)：读入 \(a\) 改变陪集，读入 \(b\) 保持陪集。选从 \(H\) 到 \(Ha\) 的那条正向 \(a\) 边为生成树，其树代表为
\[
t_H=1,\qquad t_{Ha}=a.
\]
余下正向边分别是 \(Ha\) 到 \(H\) 的另一条 \(a\) 边，以及两个顶点上的 \(b\) 环。

\begin{figure}[htbp]
\centering
\begin{tikzpicture}
\node[circle,draw,inner sep=4pt] (h) at (0,0) {\(H\)};
\node[circle,draw,inner sep=4pt] (ha) at (4,0) {\(Ha\)};
\draw[->,thick] (h) to[bend left=25] node[above] {\(a\)（树边）} (ha);
\draw[->] (ha) to[bend left=25] node[below] {\(a\)} (h);
\path[->] (h) edge[loop left] node[left] {\(b\)} (h);
\path[->] (ha) edge[loop right] node[right] {\(b\)} (ha);
\end{tikzpicture}
\caption[指数二子群的 Schreier 图]{指数二子群的 Schreier 图。只画指定的正向边，粗线表示所选树边。每条边还有未画出的反向，其标签为正向标签的逆。}
\label{b1:fig:schreier-index-two}
\end{figure}

两条方向相反却都标记为正向 \(a\) 的边，是不同的几何边。上方树边的反向虽也从 \(Ha\) 到 \(H\)，标签却是 \(a^{-1}\)，与下方正向 \(a\) 边不同。因此读入 \(aa\) 得到两条不同几何边组成的闭路径，不能作为一次往返消去；读入 \(aa^{-1}\) 才沿同一条树边返回。

\ProseParagraphBreak
逐条将非树正向边代入\eqref{b1:eq:schreier-basis}，得到\cref{b1:tab:schreier-index-two-basis}。表中第三条基环路先沿树边到达 \(Ha\)，绕 \(b\) 环后再沿树边反向返回，所以标签为 \(aba^{-1}\)。

\begin{table}[htbp]
\centering
\caption{指数二子群的非树边与自由基}
\label{b1:tab:schreier-index-two-basis}
\begin{tabular}{@{}cccc@{}}
\toprule
非树正向边 & 起点树代表 & 终点树代表 & 基元素 \(t_vs\,t_w^{-1}\) \\
\midrule
\(Ha\xrightarrow{a}H\) & \(a\) & \(1\) & \(a^2\) \\
\(H\xrightarrow{b}H\) & \(1\) & \(1\) & \(b\) \\
\(Ha\xrightarrow{b}Ha\) & \(a\) & \(a\) & \(aba^{-1}\) \\
\bottomrule
\end{tabular}
\end{table}

因此所得自由基为
\[
a^2,\qquad b,\qquad aba^{-1}.
\]
公式也给出 \(\operatorname{rank}H=3\)。它表明非交换自由群子群的秩可以大于母群的秩，不能用向量空间子空间的维数直觉判断。

\subsection{字问题与 Cayley 图初步}
\label{b1:subsec:0804:03}

给定群的有限生成集及相应展示，\term[zi-wenti]{字问题}[word problem]要求判定输入字是否表示单位。这里有限的是生成集，关系集未必有限。\cref{b1:def:free-group-construction}对任意集合 \(S\) 都给出集合论意义下的自由群，并不要求 \(S\) 具有可计算编码；若要把字的约化实现为算法，则每个输入字母必须有有限编码，还须能有效比较两个字母是否互逆。对固定的有限字母表，这些操作可视为常数时间。此时\cref{b1:lem:unique-reduced-word}给出的算法读入每个字母时至多进栈或出栈一次，而所得约化字为空恰好表示原字为单位，因而在通常按字母计费的模型下，可用与字长成正比的步骤判定。

\ProseParagraphBreak
对二面体群，\cref{b1:prop:dihedral-presentation}的证明把字整理为唯一的 \(r^is^\varepsilon\)，其中 \(0\le i<n\)、\(\varepsilon\in\{0,1\}\)，故只须检查两指数是否均为零。四面体旋转群则可依\cref{b1:eg:tetrahedral-presentation}的具体十二元素展示，先读入各字母对应的置换再逐次相乘；最终置换是否恒等即可判定原字是否为单位。

\ProseParagraphBreak
在一般展示中，不同的自由约化字仍可能表示同一个商群元素。例如在 \(\langle a\mid a^2=1\rangle\) 中，\(a\) 与 \(a^3\) 都是自由群中的约化字，却表示同一个商群元素。自由约化只处理相邻逆字母的消去，额外关系还会产生新的等式，因而它不再直接给出商群的唯一标准形。

\ProseParagraphBreak
这并不妨碍在集合意义上选取唯一代表。对有限个生成字母及其形式逆字母固定一个次序；在一个群元素的全部代表字中，先选长度最短的，再在同长的有限个字中选字典序最小的，就唯一确定了一个代表。然而，要算出这个代表，还须判断哪些候选字表示同一元素，以上选择规则本身没有提供这种判定方法。

\ProseParagraphBreak
设 \(P=\langle S\mid R\rangle\)，并取 \(w\in F(S)\)。正规闭包的有限乘积描述给出精确的等价条件：
\[
\begin{aligned}
w=1\text{ 在 }P\text{ 中}
\quad\Longleftrightarrow\quad
&w=\prod_{i=1}^{m}u_i r_i^{\varepsilon_i}u_i^{-1}
  \text{ 在 }F(S)\text{ 中成立},\\
&\text{其中 }m\ge0,\quad
  u_i\in F(S),\quad r_i\in R,\quad
  \varepsilon_i\in\{1,-1\}.
\end{aligned}
\]
这里 \(m=0\) 时右边为空乘积。因而每个真等式都有一个只含有限多个共轭关系字的证明证书。

\ProseParagraphBreak
若 \(S,R\) 都有限，就可按包括因子个数在内的总编码长度递增，枚举所有这样的有限乘积，并用自由约化检验它是否等于输入字 \(w\)。当 \(w\) 确实表示单位时，完备枚举最终会找到一个证书并终止；当 \(w\) 不表示单位时，搜索却可能永远继续。缺少的不是肯定情形的有限证书，而是对否定情形也保证终止的判定方法。本节的算法结论只针对已经给出有效编码和判定步骤的具体情形。

\ProseParagraphBreak
设 \(F(S)\) 为以 \(S\) 为基的自由群。在 \(H=\{1\}\) 时，\cref{b1:thm:nielsen-schreier}证明中构造的 Schreier 图称为 \(F(S)\) 关于 \(S\) 的\term[cayley-tu]{Cayley 图}[Cayley graph]：顶点为群元素，边记录右乘生成元。自由群的 Cayley 图是一棵树：图连通，而若有非空无往返闭路径，其标签由每个顶点的唯一字母出边性质构成非空约化字；闭路径又要求标签表示单位，这与\cref{b1:lem:unique-reduced-word}矛盾。给每条边长度一以后，从单位到一个元素的距离恰为其约化字长。于是代数乘法给出的字与图上的路径长度发生联系，这构成几何群论研究的一个基本出发点。继续学习复形、覆盖映射与 Cayley 复形，可读 Lyndon 与 Schupp 的《Combinatorial Group Theory》第 III 章第 2--4 节\cite[Chapter III, \S\S~2--4, pp. 115--124]{LyndonSchupp2001}；若要直接对照自由积，可参阅该书第 IV 章第 1 节\cite[Chapter IV, \S~1]{LyndonSchupp2001}。


\begin{chaptersummary}
自由群先由字的等价类构造，约化字唯一性再给出标准表示；当字母具有可判定的编码和比较方式时，这一过程也给出实际计算。展示通过关系字的正规闭包定义商群；验证一个具体展示，既要检查关系与生成性，也要证明所给关系没有留下多余的群元素。自由积允许独立给定两个因子的同态，直积则额外要求两像逐元素交换。

\ProseParagraphBreak
\secref{b1:sec:0804}把自由群的子群识别为 Schreier 图的闭路径群。选定生成树后，每条非树边给出一个自由基元素；有限指数时，边数与顶点数决定基的个数。下一章转入交换群，将生成与关系写成整数线性组合和关系矩阵。
\end{chaptersummary}

\BookExercises

\begin{exercise}\label{b1:ex:08-torsion-free}
证明自由群的非单位元都具有无限阶。提示：长度至多为一或首尾字母不互逆的约化字称为\term[xunhuan-yuehua-zi]{循环约化字}[cyclically reduced word]。先把一个非空约化字写成 \(uvu^{-1}\)，其中 \(v\) 非空且循环约化。
\end{exercise}
\begin{solution}
若约化字首尾互逆，就把这两个字母分别移入 \(u\) 及 \(u^{-1}\)，对中间字重复。长度每次减少二；非空约化字不能最终只剩一对相邻逆字母，所以最后得到非空且循环约化的 \(v\)。此时 \(v^m\) 对每个正整数 \(m\) 仍约化且非空。原字是 \(uvu^{-1}\) 的约化拼接，故 \(uv^mu^{-1}\) 也在相同的两个接缝处不消去。由\cref{b1:lem:unique-reduced-word}，非空约化字不能表示单位，因此 \(uv^mu^{-1}\ne1\)。

乘方时，中间的 \(u^{-1}u\) 相消。例如
\[
(uvu^{-1})^2
=uv\underbrace{u^{-1}u}_{1}vu^{-1}
=uv^2u^{-1}.
\]
逐次使用这一消去，归纳得到 \((uvu^{-1})^m=uv^mu^{-1}\)。所以原字的任何正幂都非单位。
\end{solution}

\begin{exercise}\label{b1:ex:08-squares-embedding}
在 \(F=F\bigl(\{a,b\}\bigr)\) 中令 \(H=\langle a^2,b^2\rangle\)。
\begin{enumerate}
\item 证明 \(a^2,b^2\) 自由生成真子群 \(H<F\)。
\item 利用到 \(C_2*C_2\) 的满同态证明 \([F:H]=\infty\)。
\end{enumerate}
\end{exercise}
\begin{solution}
1. 由\cref{b1:thm:free-group-universal}，存在唯一同态 \(F\bigl(\{x,y\}\bigr)\to F\bigl(\{a,b\}\bigr)\) 发送 \(x\mapsto a^2,y\mapsto b^2\)。一个非空约化字可按相同生成元的最大块写成 \(x^{m_1}y^{n_1}\cdots\) 一类交替表达式，所有出现的指数非零；端部可以从任一类型开始或结束。替换后各指数乘二，交替的非零块之间没有消去。由\cref{b1:lem:unique-reduced-word}，像非单位，故同态的核只有单位元；它因而单射，并给出到 \(\langle a^2,b^2\rangle\) 的同构。为检验真包含，把 \(C_2=\mathbb Z/2\mathbb Z\) 写成加法群，令 \(a\mapsto1,b\mapsto0\)，并由自由群的泛性质延拓为同态。这里 \(0\) 是单位元，\(1\) 是非零元；\(a^2,b^2\) 的像均为零，而 \(a\) 的像非零。因此生成子群不含 \(a\)，是真子群。

2. 设 \(x,y\) 分别为 \(C_2*C_2\) 的两个二阶因子的非单位元。由\cref{b1:thm:free-group-universal}，映射 \(a\mapsto x,b\mapsto y\) 延拓为满同态
\[
\pi:F\longrightarrow C_2*C_2.
\]
因为 \(x^2=y^2=1\)，有 \(H=\langle a^2,b^2\rangle\subseteq\ker\pi\)。对任意 \(h\in H\)，
\[
\pi(hg)=\pi(h)\pi(g)=\pi(g),
\]
所以 \(\pi\) 在每个右陪集 \(Hg\) 上取同一值。每个右陪集只能贡献一个像元素；若 \([F:H]\) 有限，则满射 \(\pi\) 的像也只能有有限多个元素。但\secref{b1:sec:0803}已由交错约化字证明 \(C_2*C_2\) 无限，矛盾。这里仅用右陪集计数，不要求 \(H\) 正规。

也可用 Schreier 秩公式证明：第一问说明 \(H\) 是秩为二的自由群。若其指数是有限数 \(d\)，\cref{b1:thm:schreier-rank}给出
\[
2=\operatorname{rank}H=1+d(2-1)=1+d,
\]
所以 \(d=1\)。这会推出 \(H=F\)，与第一问证明的真包含矛盾。因此 \([F:H]=\infty\)。这也说明无限指数子群仍可能具有有限秩。
\end{solution}

\begin{exercise}\label{b1:ex:08-quaternion-presentation}
证明
\[
Q_8\cong\langle x,y\mid x^4=1,\ y^2=x^2,\ yxy^{-1}=x^{-1}\rangle.
\]
\end{exercise}
\begin{solution}
按\eqref{b1:eq:quaternion-relations}及其相邻的反向乘法规则，\(i^2=j^2=-1\)、\(ji=-ij\)。因此 \(i^4=1\)、\(j^2=i^2\)，且 \(jij^{-1}=-i=i^{-1}\)。把 \(x,y\) 分别映为 \(i,j\)，由\cref{b1:prop:presentation-homomorphism}得到展示群到 \(Q_8\) 的同态；其像含 \(i,j\)，从而含 \(\pm1,\pm i,\pm j,\pm ij\)，故满射。

在展示群中，先由 \(x^4=1\)、\(y^4=x^4=1\) 把负指数改成非负指数；再用 \(yx=x^{-1}y\) 把 \(y\) 推到右端。随后用 \(y^2=x^2\) 把 \(y\) 的指数减至零或一，再由 \(x^4=1\) 把 \(x\) 的指数减至 \(0,1,2,3\)。

现在比较元素个数。每个元素都可写为 \(x^iy^\varepsilon\)，所以展示群至多八个元素。已有满同态到八阶群，它也至少有八个元素；因此该满同态为同构。

关系 \(y^2=x^2\) 不能改成 \(y^2=1\)；后者正是\cref{b1:prop:dihedral-presentation}中 \(n=4\) 的展示，给出 \(D_8\)。
\end{solution}

\begin{exercise}\label{b1:ex:08-two-reflection-presentation}
对 \(n\ge3\)，证明
\[
\bigl\langle a,b\mid a^2=b^2=(ab)^n=1\bigr\rangle\cong D_{2n}.
\]
\end{exercise}
\begin{solution}
在左边令 \(r=ab,s=a\)。有 \(r^n=s^2=1\)，并且
\(srs^{-1}=ba=r^{-1}\)，由\cref{b1:prop:dihedral-presentation}，\(D_{2n}\) 正由这三个关系展示，再用\cref{b1:prop:presentation-homomorphism}，得到 \(D_{2n}\) 到左边展示群的同态。反向在 \(D_{2n}\) 中令 \(a=s,b=sr\)，则 \(a^2=s^2=1\)、\(b^2=srsr=r^{-1}r=1\)、\(ab=r\)，所以也满足 \((ab)^n=1\)；同一展示的映射性质给出反向同态。两方向在相应生成元上复合为恒等，故互为逆同构。这个展示强调两个反射，正文的展示强调旋转与反射。
\end{solution}

\begin{exercise}\label{b1:ex:08-free-product-abelianization}
对群 \(G\) 记 \(G_{\mathrm{ab}}=G/G'\)。证明
\[
(G*H)_{\mathrm{ab}}\cong G_{\mathrm{ab}}\times H_{\mathrm{ab}}.
\]
\end{exercise}
\begin{solution}
记 \(F=G*H\)、\(B=F/F'\)。将 \(g\in G,h\in H\) 分别送到 \((gG',1)\)、\((1,hH')\)，由\cref{b1:thm:free-product-universal}得到满同态 \(u:F\to G_{\mathrm{ab}}\times H_{\mathrm{ab}}\)。目标群是交换群，故
\[
u\bigl([x,y]\bigr)=\bigl[u(x),u(y)\bigr]=1,
\]
从而 \(F'\subseteq\ker u\)。由\cref{b1:thm:quotient-universal}，它唯一地诱导 \(\bar u:B\to G_{\mathrm{ab}}\times H_{\mathrm{ab}}\)。

反向把因子 \(G,H\) 映入 \(F\) 后再商到 \(B\)。由\cref{b1:prop:derived-quotient}，\(B\) 交换；上述交换子计算说明这两个映射分别消去 \(G',H'\)，所以同一商群泛性质给出 \(\alpha:G/G'\to B\)、\(\beta:H/H'\to B\)。定义
\[
v(gG',hH')=\alpha(gG')\beta(hH').
\]
两像在 \(B\) 中交换，所以 \(v\) 保持乘法。复合 \(\bar uv\) 固定直积的两个因子，\(v\bar u\) 固定 \(B\) 中 \(G,H\) 的像；这些像分别生成相应群，因此两复合都是恒等。
\end{solution}

\begin{exercise}\label{b1:ex:08-redundant-generator}
设 \(P=\langle S\mid R\rangle\)，\(w\in F(S)\)。
\begin{enumerate}
\item 取新符号 \(t\notin S\)。证明
\[
Q=\bigl\langle S\sqcup\{t\}\mid R,\ t=w\bigr\rangle
\]
与 \(P\) 同构，并说明为何只增加 \(t\) 而不增加关系通常会改变群。
\item 证明自然满同态
\[
\rho:\langle S\mid R\rangle
\longrightarrow
\langle S\mid R\cup\{w\}\rangle
\]
为同构，当且仅当
\[
w\in\NormalClosure{R}_{F(S)}.
\]
\end{enumerate}
\end{exercise}
\begin{solution}
1. 映射 \(S\to P\) 取原有生成元在 \(P\) 中的像，并令 \(t\mapsto\overline w\)，满足 \(Q\) 的全部关系，故\cref{b1:prop:presentation-homomorphism}给出同态 \(Q\to P\)。由 \(S\) 在 \(Q\) 中仍满足 \(R\)，同一命题也给出 \(P\to Q\)。复合在 \(S\) 上均为恒等，在 \(Q\) 的 \(t\) 上则由 \(t=w\) 得到恒等，因此两者互逆。若只加新符号，在自由积的展示公式\eqref{b1:eq:free-product-presentation}中取第二个群为 \(\mathbb Z=\langle t\mid\varnothing\rangle\)：两组生成符号不交，合并后的关系仍只有 \(R\)，所以所得群为 \(P*\mathbb Z\)。例如 \(P=1\) 时它变为 \(\mathbb Z\)，所以缺少关系会引入新的自由选择。

2. 在 \(F(S)\) 中记
\[
N=\NormalClosure{R}_{F(S)},\qquad
M=\NormalClosure{R\cup\{w\}}_{F(S)}.
\]
因为 \(N\leq M\)，商映射诱导自然满同态 \(\rho:F(S)/N\rightarrow F(S)/M\)，其核为 \(M/N\)。因此 \(\rho\) 为同构，当且仅当 \(M=N\)。若 \(M=N\)，则 \(w\in M=N\)；反过来，若 \(w\in N\)，则正规子群 \(N\) 已包含 \(R\cup\{w\}\)，故 \(M\leq N\)，再结合 \(N\leq M\) 得 \(M=N\)。这就证明了所述充要条件。
\end{solution}

\begin{optionalexercise}\label{b1:ex:08-index-three}
设 \(F=F\bigl(\{a,b\}\bigr)\)，把 \(C_3=\mathbb Z/3\mathbb Z\) 写成加法群，令 \(H\) 为同态 \(F\to C_3\)、\(a\mapsto1,b\mapsto0\) 的核。画出 Schreier 图，选从 \(H\) 到 \(Ha\) 及从 \(Ha\) 到 \(Ha^2\) 的两条正向 \(a\) 边为生成树，树代表为 \(1,a,a^2\)。证明相应的自由基为
\[
A=a^3,\qquad B_0=b,\qquad
B_1=aba^{-1},\qquad B_2=a^2ba^{-2}.
\]
再沿 Schreier 图读取 \(w=ababab^{-1}\)，把它改写成 \(A,B_0,B_1,B_2\) 的字，并说明原自由群约化与新自由基约化所使用的字母表有何不同。提示：在新字母表中已约化的字，展开为 \(a,b\) 及其逆字母后，接缝处仍可能发生消去。
\end{optionalexercise}
\begin{solution}
陪集为 \(H,Ha,Ha^2\)。三条正向 \(a\) 边依次构成三角形，每个顶点另有一个 \(b\) 环。取前两条 \(a\) 边为树，树代表字为 \(1,a,a^2\)。非树 \(a\) 边给出 \(a^3\)，三个 \(b\) 环分别给出
\[
b,\qquad aba^{-1},\qquad a^2ba^{-2}.
\]
具体地，按\eqref{b1:eq:schreier-basis}，末条 \(a\) 边给出 \(a^2a\cdot1^{-1}=a^3\)，各顶点的 \(b\) 环给出 \(a^kba^{-k}\)（\(k=0,1,2\)）。\cref{b1:lem:graph-loop-free}断言每条非树边恰对应一个自由基元素，而式\eqref{b1:eq:schreier-label-isomorphism}把闭路径群同构地识别为 \(H\)，所以这四个元素确实构成自由基。由\cref{b1:thm:first-isomorphism}，该满同态的核指数为三；\cref{b1:thm:schreier-rank}也给出秩 \(1+3(2-1)=4\)。

现在沿图读取
\[
w=ababab^{-1}.
\]
第一次 \(b\) 出现在顶点 \(Ha\)，记录 \(B_1\)；第二次 \(b\) 出现在顶点 \(Ha^2\)，记录 \(B_2\)；随后从 \(Ha^2\) 沿非树 \(a\) 边回到根，记录 \(A\)；最后在根处反向绕 \(b\) 环，记录 \(B_0^{-1}\)。因此
\[
w=B_1B_2AB_0^{-1}.
\]
在原基中展开右端得到
\[
(aba^{-1})(a^2ba^{-2})a^3b^{-1}
=ababab^{-1},
\]
与原字一致。表达式 \(B_1B_2AB_0^{-1}\) 作为新自由基字母上的字已经约化；把各个基元素展开为原字母 \(a,b\) 后，还须在接缝处进行自由消去，才得到右端的原字 \(ababab^{-1}\)。两种约化分别使用子群的新自由基字母表和母群的原自由基字母表。
\end{solution}

\begin{optionalexercise}\label{b1:ex:08-infinite-rank-kernel}
设 \(F=F\bigl(\{a,b\}\bigr)\)，同态 \(F\to\mathbb Z\) 发送 \(a\mapsto1,b\mapsto0\)。记该同态的核为 \(H\)。证明 \(H\) 自由，且自由基为 \(\{a^kba^{-k}\mid k\in\mathbb Z\}\)；再证明 \(H=\NormalClosure{b}_{F}\)，并比较正规生成与普通子群生成。
\end{optionalexercise}
\begin{solution}
核的右陪集由整数 \(k\) 标记，代表字取 \(a^k\)。正向 \(a\) 边从 \(k\) 指向 \(k+1\)，正向 \(b\) 边是各顶点的环。所有 \(a\) 边组成一条双向无限的树，所有非树边正好是 \(b\) 环。由\eqref{b1:eq:schreier-basis}，位于 \(k\) 的 \(b\) 环对应 \(a^kba^{-k}\)。\cref{b1:lem:graph-loop-free}与式\eqref{b1:eq:schreier-label-isomorphism}说明，所有这些非树边对应的元素构成 \(H\) 的自由基。若 \(H\) 有有限生成集，这些生成元的约化字总共只涉及有限个基字母，其乘积不可能产生其余基字母，矛盾。因此该核不是有限生成群，尽管母群只需两个生成元。

令 \(N=\NormalClosure{b}_{F}\)。因 \(b\in H\) 且 \(H\) 是同态的核，\cref{b1:def:normal-closure}给出 \(N\le H\)。反过来，正规性使 \(N\) 包含每个 \(a^kba^{-k}\)，而这些元素已证为 \(H\) 的自由基，所以 \(H\le N\)。于是 \(H=N\)。这个例子把两种生成方式分开：在母群 \(F\) 中，只需 \(b\) 一个正规生成元；在 \(H\) 自身中，普通子群生成需要无限多个元素。因而有限正规生成不蕴含有限普通生成。这里正规闭包取在 \(F\) 中，而不是取在 \(H\) 自身中。
\end{solution}

\begin{optionalexercise}\label{b1:ex:08-nonnormal-schreier}
设 \(F=F\bigl(\{a,b\}\bigr)\)，令同态 \(\varphi:F\to S_3\) 满足
\(\varphi(a)=(12)\)、\(\varphi(b)=(123)\)。取自然置换作用中点 \(1\) 的稳定子群
\[
L=\operatorname{Stab}_{S_3}(1)=\{1,(23)\},\qquad K=\varphi^{-1}(L).
\]
证明 \(K\) 的指数为三，但不正规；构造它的 Schreier 图，选一棵生成树并求自由基。
\end{optionalexercise}
\begin{solution}
\cref{b1:thm:free-group-universal}保证所述同态存在。其像含二阶元和三阶元，由\cref{b1:thm:lagrange}，像的阶同时被 \(2,3\) 整除，所以像为整个 \(S_3\)。\(L\) 是稳定子群，其原像 \(K\) 对乘法与逆封闭，故为 \(F\) 的子群。

同一右陪集的代表由左乘 \(k\in K\) 相互转换，而 \(\varphi(k)\) 固定点 \(1\)。因此用逆置换记录点的位置，就有
\[
\varphi(kg)^{-1}(1)
=\varphi(g)^{-1}\varphi(k)^{-1}(1)
=\varphi(g)^{-1}(1).
\]
这说明逆号使标记在整个右陪集上保持一致。它也使右端乘入 \(s\) 对应在标记上左作用 \(\varphi(s)^{-1}\)，因为 \(\varphi(gs)^{-1}=\varphi(s)^{-1}\varphi(g)^{-1}\)。据此定义
\[
\Psi(Kg)=\varphi(g)^{-1}(1).
\]
若 \(Kg=Kh\)，则 \(h=kg\) 对某个 \(k\in K\) 成立；由于 \(\varphi(k)^{-1}\) 固定 \(1\)，两代表得到的标记相同。反过来，若 \(\varphi(g)^{-1}(1)=\varphi(h)^{-1}(1)\)，在两边作用 \(\varphi(h)\)，便知 \(\varphi(hg^{-1})\) 固定 \(1\)，即 \(hg^{-1}\in K\)，从而 \(Kh=Kg\)。代表字 \(1,a,b\) 分别给出标记 \(1,2,3\)，因此全部右陪集恰为 \(K,Ka,Kb\)，指数为三。

对正字母 \(s\in\{a,b\}\)，沿边右乘满足
\[
\Psi(Kgs)=\varphi(s)^{-1}\bigl(\Psi(Kg)\bigr).
\]
于是图的全部正向边为
\[
\begin{aligned}
a:&\quad K\longrightarrow Ka,\quad Ka\longrightarrow K,\quad Kb\longrightarrow Kb;\\
b:&\quad K\longrightarrow Kb,\quad Ka\longrightarrow K,\quad Kb\longrightarrow Ka.
\end{aligned}
\]
每个箭头代表一条几何边的指定正向，反向另标相应逆字母。选 \(K\xrightarrow{a}Ka\)、\(K\xrightarrow{b}Kb\) 为生成树，树代表为 \(t_K=1,t_{Ka}=a,t_{Kb}=b\)。其余四条正向边代入\eqref{b1:eq:schreier-basis}，依次得到
\[
\begin{aligned}
Ka\xrightarrow{a}K&:\quad a\cdot a\cdot1^{-1}=a^2,\\
Kb\xrightarrow{a}Kb&:\quad b\cdot a\cdot b^{-1}=bab^{-1},\\
Ka\xrightarrow{b}K&:\quad a\cdot b\cdot1^{-1}=ab,\\
Kb\xrightarrow{b}Ka&:\quad b\cdot b\cdot a^{-1}=b^2a^{-1}.
\end{aligned}
\]
\cref{b1:lem:graph-loop-free}给出非树边对应的闭路径自由基，而式\eqref{b1:eq:schreier-label-isomorphism}将它同构地送到 \(K\)，故这四个元素是 \(K\) 的自由基。\cref{b1:thm:schreier-rank}也给出秩 \(1+3(2-1)=4\)。

最后，按从右向左作用的置换乘法，\(\varphi(ab)=(23)\in L\)，而
\[
\varphi\bigl(a(ab)a^{-1}\bigr)=(12)(23)(12)=(13)\notin L.
\]
所以 \(ab\in K\)，但 \(a(ab)a^{-1}\notin K\)，\(K\) 不正规。这张图由右陪集构成，并不需要把这些陪集定义成商群。
\end{solution}
